% ============================================================================
\chapter{Daily operation}
\label{ch:usage}
% ============================================================================

This chapter describes operation through the touch screen only; no keyboard is
required. The navigation between screens is summarized in the screen map of
Section~\ref{sec:screenmap}.

\section{Start-up and the status bar}

After power-on (with the systemd unit of Section~\ref{sec:systemd}), the application
starts automatically and immediately shows the \textbf{Main} screen while it connects
to the robot in the background. The status bar at the bottom shows:

\begin{itemize}[itemsep=2pt]
  \item a \textbf{connection indicator}: green \emph{Connected} once the robot has
        answered the identification handshake, red \emph{Disconnected} otherwise;
  \item the latest \textbf{status message}, e.g.\ \emph{Connecting\dots},
        \emph{Connected}, \emph{Move OK}, \emph{Homing complete}, or an error text.
\end{itemize}

\begin{note}
If the robot was off when the terminal started, the status shows
\emph{Connection \dots\ failed}. Power the robot on, then restart the application
(power-cycle the terminal, or \code{sudo systemctl restart deltax2} over SSH).
An in-UI reconnect button is planned (Appendix~\ref{app:limitations}).
\end{note}

\section{Main screen}

Three buttons:

\begin{center}
\begin{tabularx}{\textwidth}{@{}lX@{}}
  \toprule
  \textbf{Button} & \textbf{Function} \\
  \midrule
  \textbf{Seeding} & Starts the guided seeding workflow (Section~\ref{sec:seeding}). \\
  \textbf{Calibration} & Opens manual jog and homing (Section~\ref{sec:calibration}). \\
  \textbf{Configuration} & Opens system utilities (Section~\ref{sec:systemscreen}). \\
  \bottomrule
\end{tabularx}
\end{center}

\section{Seeding a plate}
\label{sec:seeding}

\begin{enumerate}[itemsep=3pt]
  \item Touch \textbf{Seeding} on the main screen.
  \item \textbf{Install the seeding plate}: place the tray in its mechanical
        reference position, then touch the button matching its type
        (e.g.\ \textbf{77pots}). The choice is recorded and shown in the status bar.
        \textbf{Cancel} returns to the main screen.
  \item \textbf{Load the seeds}: fill the seed supply, then touch \textbf{Ok} to
        start the job --- the robot homes first and then processes the tray pot by
        pot. \textbf{Cancel} returns to the main screen without moving the robot.
  \item \textbf{Seeding in progress}: the screen shows the live progress
        (\emph{Pot 12 / 77}) and three controls:
    \begin{itemize}[itemsep=1pt]
      \item \textbf{Pause} ($\parallel$) --- the robot finishes the current pot and
            waits. The status bar shows \emph{Seeding paused}.
      \item \textbf{Continue} ($\triangleright$) --- resumes a paused job.
      \item \textbf{Stop} ($\square$) --- aborts the job: the robot finishes the
            current pot, then stops. The display returns to the main screen and the
            status bar shows \emph{Stopping\dots} followed by \emph{Seeding stopped}.
    \end{itemize}
  \item On completion the display returns to the main screen and the status bar
        shows \emph{Seeding complete}.
\end{enumerate}

\begin{warning}
\textbf{Stop} waits for the current pot to finish; it does not cut power. For
immediate mechanical arrest in an emergency, use the machine's hardware
emergency stop (if fitted) or remove robot power.
\end{warning}

\section{Calibration and manual movement}
\label{sec:calibration}

The \textbf{Calibration} screen provides:

\begin{itemize}[itemsep=2pt]
  \item \textbf{Step size selector} --- four buttons: 0.01, 0.1, 1 or 5\,mm.
        The selected value applies to every subsequent jog.
  \item \textbf{XY pad} --- arrow buttons jog the head in X and Y by one step per
        touch; the centre button homes all axes (\code{G28}).
  \item \textbf{Z column} --- up/down jog the head vertically; the centre button
        also triggers homing.
  \item \textbf{Cart column} --- up/down adjust the rotation axis; the centre
        button resets its displayed value (see the note below).
  \item \textbf{Position panel} --- live display of the logical X, Y, Z and Cart
        coordinates in mm.
\end{itemize}

Each jog is checked against the safety limits; a refused move leaves the robot
untouched and shows \emph{Movement out of safety limits} in the status bar.
After homing, the position display resets to $0.000$ on all axes.

\begin{note}
In the current software version the \textbf{Cart} axis is tracked on screen but not
yet driven on the hardware (Appendix~\ref{app:limitations}).
\end{note}

\begin{tip}
Jog commands are queued and executed in order. On a touch screen it is easy to tap
faster than the robot moves --- prefer single deliberate taps and watch the position
readout; each completed move updates the display.
\end{tip}

\section{System screen}
\label{sec:systemscreen}

The \textbf{Configuration} button opens the system screen, which currently offers:

\begin{itemize}[itemsep=2pt]
  \item \textbf{List com} --- enumerates the serial ports detected on the system.
        The result is written to the application log (visible with
        \code{journalctl -u deltax2} over SSH), which helps an administrator find
        the correct \code{port} value for \code{config.toml}.
\end{itemize}
